\subsection{Structure of dualizing modules}

Lemma 2.--- Let A be a Noetherian ring, and let M and N be finitely generated A-modules, \(u: \mathrm{M} \to \mathrm{N}\) a homomorphism. Let \(x\) an element of the radical of A, such that the homothety \(x_{\mathrm{N}}\) be injective. If the homomorphism \(\overline{u}: \mathrm{M}/x\mathrm{M} \to \mathrm{N}/x\mathrm{N}\) induced by \(u\) is injective (resp. surjective, resp. bijective), then so is \(u\) .

The assertion concerning the surjectivity of \(u\) follows from Nakayama's lemma (II, § 3, no. 2, cor. 1 of prop. 4), without any hypothesis on \(x_{\mathrm{N}}\) . Consider the commutative diagram with exact rows

\[
\begin{array}{c c c c c c c} & \mathrm{M} & \xrightarrow {x _ {\mathrm{M}}} & \mathrm{M} & \longrightarrow & \mathrm{M} / x \mathrm{M} \\ & u \Bigg \downarrow & & u \Bigg \downarrow & & \overline {{u}} \Bigg \downarrow \\ 0 & \longrightarrow & \mathrm{N} & \xrightarrow {x _ {\mathrm{N}}} & \mathrm{N} & \longrightarrow & \mathrm{N} / x \mathrm{N} \end{array} ;
\]

using the snake lemma (I, § 1, n° 4, prop. 2), one deduces an exact sequence \(\operatorname{Ker} u \xrightarrow{x} \operatorname{Ker} u \longrightarrow \operatorname{Ker} \overline{u}\) If \(\overline{u}\) is injective, the homothety with ratio \(x\) is surjective in \(\operatorname{Ker} u\) , which implies \(\operatorname{Ker} u = 0\) by Nakayama's lemma.

PROPOSITION 7.--- Let A be a Noetherian ring and Ω a dualizing A-module.

\begin{enumerate}
\def\labelenumi{\alph{enumi})}
\item
  We have \(\operatorname{Ext}_{\mathrm{A}}^{i}(\Omega, \Omega) = 0\) for \(i > 0\) .
\item
  \(L'\) canonical homomorphism \(\gamma : \mathrm{A} \to \operatorname{End}_{\mathrm{A}}(\Omega)\) is bijective.
\item
  Every dualizing A-module is of the form \(\Omega\otimes_{\mathrm{A}}\mathrm{L}\) where \(\mathrm{L}\) is a projective A-module of rank 1.
\end{enumerate}

\begin{enumerate}
\def\labelenumi{\Alph{enumi})}
\tightlist
\item
  Let us first treat the case where the ring A is local. In this case, condition c) simply means that two dualizing modules are isomorphic.
\end{enumerate}

Let \(\Omega'\) a dualizing A-module. We have \(\operatorname{prof}_{\mathrm{A}}(\Omega') = \dim_{\mathrm{A}}(\Omega') = \dim(\mathrm{A})\) (no. 1, prop. 1), hence \(\operatorname{Ext}_{\Lambda}^{i}(\Omega', \Omega) = 0\) for \(i \neq 0\) (no. 1, prop. 3, c)), whence a).

Let us prove b) and c) by induction on the integer dim(A) (equal to prof(A)). If it is zero, the ring A is Artinian, \(\Omega'\) and \(\Omega\) are Matlis A-modules ( \(n^{\circ}\) 1, example 1); they are therefore isomorphic ( \(\S\) 8, \(n^{\circ}\) the canonical map A → \(\mathrm{End}_{\mathrm{A}}(\Omega)\) is bijective \(\S\) 8, \(n^{\circ}\) 2, prop. 3, c)). Suppose dim(A) \textgreater{} 0 and let \(x\) be a cancellable element of \(\mathfrak{m}_{\mathrm{A}}\) . The homothety \(x_{\Omega}\) is injective ( \(n^{\circ}\) 2, prop. 4), and one has an exact sequence

\[
0 \rightarrow \Omega \xrightarrow {x _ {\Omega}} \Omega \longrightarrow \Omega / x \Omega \rightarrow 0.
\]

Since \(\mathrm{Ext}_{\mathrm{A}}^{1}(\Omega',\Omega)\) is zero and that \(\mathrm{Hom}_{\mathrm{A}}(\Omega',\Omega /x\Omega)\) is identified with \(\mathrm{Hom}_{\mathrm{A} / x\mathrm{A}}(\Omega'/x\Omega',\Omega/x\Omega)\) , one deduces an exact sequence

\[
0 \rightarrow \operatorname{Hom} _ {\mathrm{A}} \left(\Omega^ {\prime}, \Omega\right) \xrightarrow {x} \operatorname{Hom} _ {\mathrm{A}} \left(\Omega^ {\prime}, \Omega\right) \xrightarrow {p} \operatorname{Hom} _ {\mathrm{A} / x \mathrm{A}} \left(\Omega^ {\prime} / x \Omega^ {\prime}, \Omega / x \Omega\right)\rightarrow 0,\tag{1}
\]

where \(p\) is the canonical map. By prop. 4, the A/xA-modules \(\Omega/x\Omega\) and \(\Omega'/x\Omega'\) are dualizing, hence isomorphic by the induction hypothesis. Let \(\overline{u}\) an isomorphism of \(\Omega'/x\Omega'\) on \(\Omega/x\Omega\) . Taking into account the exact sequence (1), there exists an A-homomorphism \(u: \Omega' \to \Omega\) such that \(p(u) = \overline{u}\) ; by Lemma 2, \(u\) is bijective, which proves c). By the induction hypothesis, the canonical homomorphism A/xA \(\longrightarrow\) End \(_{A/xA}(\Omega/x\Omega)\) is bijective. Given the exact sequence (1), this homomorphism is identified with the homomorphism \(\overline{\gamma}: A/xA \longrightarrow \text{End}_A(\Omega)/x \text{End}_A(\Omega)\) induced by \(\gamma\) it then follows from Lemma 2 that \(\gamma\) is bijective, whence b).

\begin{enumerate}
\def\labelenumi{\Alph{enumi})}
\setcounter{enumi}{1}
\tightlist
\item
  Let us turn to the general case. For every maximal ideal m of A and every integer i \textgreater{} 0, we have \(\operatorname{Ext}_{A_{m}}^{i}(\Omega_{\mathfrak{m}}, \Omega_{\mathfrak{m}}) = 0\) based on the foregoing, therefore \(\operatorname{Ext}_{A}^{i}(\Omega, \Omega)_{\mathfrak{m}} = 0\) (§ 3, no. 2, prop. 2), which implies \(\operatorname{Ext}_{A}^{i}(\Omega, \Omega) = 0\) (II, § 3, no. 3, cor. 2 of th. 1). Likewise, the homomorphism \(\gamma_{m}: A_{m} \to \operatorname{End}_{A}(\Omega)_{m}\) is bijective for every maximal ideal m of A, hence \(\gamma\) is bijective (loc. cit., th. 1).
\end{enumerate}

Let us finally prove c). Let \(\Omega'\) a dualizing A-module. Denote by L the A-module \(\mathrm{Hom}_{\mathrm{A}}(\Omega',\Omega)\) and by \(v:\Omega' \otimes_{\mathrm{A}} \mathrm{L} \longrightarrow \Omega\) the homomorphism such that \(v(x \otimes f) = f(x)\) for \(x \in \Omega'\) , \(f \in \mathrm{L}\) . Let \(\mathfrak{m}\) a maximal ideal of A. The \(\mathrm{A}_{\mathfrak{m}}\) -module \(\mathrm{L}_{\mathfrak{m}}\) is identified with \(\mathrm{Hom}_{\mathrm{A}_{\mathfrak{m}}}(\Omega_{\mathfrak{m}}',\Omega_{\mathfrak{m}})\) by the already treated case it is free of rank one, and every isomorphism \(h:\Omega_{\mathfrak{m}}' \to \Omega_{\mathfrak{m}}\) is a generator. When one identifies \(\mathrm{L}_{\mathfrak{m}}\) with \(\mathrm{A}_{\mathfrak{m}}\) using the generator \(h\) , the homomorphism \(v_{\mathfrak{m}}:\Omega_{\mathfrak{m}}' \otimes_{\mathrm{A}_{\mathfrak{m}}}\mathrm{L}_{\mathfrak{m}} \longrightarrow \Omega_{\mathfrak{m}}\) is identified with \(h\) therefore it is bijective. Since this holds for every maximal ideal \(\mathfrak{m}\) of A, the A-module L is projective of rank one (II, § 5, no. 3, th. 2), and the homomorphism \(v\) is bijective (II, § 3, no. 3, th. 1).

COROLLARY 1.--- For A to be a Gorenstein ring, it is necessary and sufficient that the A-module Ω be projective of rank 1.

This follows from example 2 of no. 1 and from prop. 7 c).

COROLLARY 2.--- Suppose that the ring A is presentable. The set of prime ideals p of A such that A \(_{p}\) let a Gorenstein ring be open in Spec(A).

Let p be a prime ideal of A such that \(A_{p}\) Let A be a Gorenstein ring. It is then a Macaulay ring; after possibly replacing A by \(A_{f}\) , where f is a suitable element of A - p, one reduces to the case where A is a Macaulay ring (§ 4, n° 4, prop. 7, c)). Let \(\Omega\) be a dualizing A-module (n° 3, cor. 3 of prop. 6). Then \(\Omega_{p}\) is a dualizing module over the Gorenstein ring \(A_{p}\) (n° 1, prop. 2), hence is free of rank 1 (cor. 1). Consequently there exists an element g of A - p such that the \(A_{g}\) -module \(\Omega_{g}\) is free of rank 1 (II, § 5, n° 1, cor. of prop. 2). Thus \(A_{g}\) is a Gorenstein ring (cor. 1) and the same is true of \(A_{q}\) for every prime ideal q of A not containing g (§ 3, n° 7, example 1), which proves the corollary.

